% ==========================================================
% 示范章节：同一份源码，四种成品（约 40 行，只演示「怎么写」）
% 挂载：main.tex 中把 \demomodefalse 改为 \demomotrue，
%       或运行 python3 build_all.py --demo
% 覆盖写法：通用内容 + advancedonly + basiconly + \vchoice{}{}
%          + \tk{} + solution + warmupbox/trapbox/example
% ==========================================================
\section{示范：四合一多版本写法}

\begin{warmupbox}{读法（先看这里）}
本节不教新知识，只演示\textbf{版本分流写法}：修改 \verb|main.tex| 顶部的两行开关
（\verb|\teachertrue/\teacherfalse| 与 \verb|\basictrue/\basicfalse|），
即可从\emph{同一份源码}导出四种成品 PDF，不必复制、粘贴或维护任何副本。
\end{warmupbox}

设 $f(x)=x^{2}-2x$，考察它在区间 $[0,2]$ 上的结构。

\begin{learnobjective}
\begin{itemize}
  \item 会算 $f(0),f(1),f(2)$ 三个值并描点画草图；
  \item 会用集合数对、几何位置、代数方程三种语言描述同一条抛物线；
  \item 初步看懂「平方算子」$f\mapsto f^{2}$ 的零点行为。
\end{itemize}
\end{learnobjective}

\subsection*{第 (1) 问：具象感知与画图}
计算：$f(0)=\tk{0}$，$f(1)=\tk{-1}$，$f(2)=\tk{0}$；
描出 $(0,0),(1,-1),(2,0)$ 三点，并画出草图。

\vchoice{%
  三个函数值对应三个点；用集合数对写即
  $\{(x,y)\mid y=x^{2}-2x,\ x\in[0,2]\}$，几何上是开口向上、对称轴为
  $x=1$ 的一段抛物线。}{%
  三个点：$(0,0)$、$(1,-1)$、$(2,0)$；连起来是开口向上的小弧线，
  最低点在 $x=1$ 处。}

\subsection*{第 (2) 问：常规通法演算}
配方：$f(x)=x^{2}-2x=\tk{(x-1)^{2}-1}$，故顶点为 $\tk{(1,-1)}$，
对称轴为直线 $\tk{x=1}$。

\subsection*{第 (3) 问：三视角结构翻译}
同一条曲线可以写成三种语言：集合语言
$\{(x,f(x))\mid x\in\mathbb{R}\}$；几何语言「以 $x=1$ 为对称轴的抛物线」；
代数方程 $y=(x-1)^{2}-1$。三者说的是\emph{同一结构}。

\begin{advancedonly}
\subsection*{第 (4) 问：高观点降维秒杀（仅完整版可见）}
作平移 $u=x-1$，则 $y=u^{2}-1$：绝对尺度被消去，只剩一个比例变量 $u$，
这就是齐次化与平移归一的思想雏形。进一步看平方算子：
$f^{2}(x)=(x(x-2))^{2}$ 的零点是\textbf{水平相切的偶重零点}
（局部 $\approx 4(x-1)^{2}$ 的二重根平移），不是尖点——它解释了
「平方后符号不变、零点处压平」的现象。
\end{advancedonly}

\begin{basiconly}
\subsection*{第 (4)' 问：大白话口诀与半成品填空（仅薄弱生版可见）}
口诀：「平方不改号，零点压平地；倒数撕裂缝，无穷变成零。」
照样子填空：$f(1)=1-2=\tk{-1}$；顶点横坐标 $x=\tk{1}$。
\end{basiconly}

\begin{example}[用两种语言看同一结构]
\vchoice{双射同构映射}{一一对应的加工机器}的说法在本讲统一替换为
\vchoice{二次型结构}{抛物线的对称结构}。请用「顶点式」写出
$f(x)=x^{2}-4x+3$ 的对称轴：$\tk{x=2}$。
\end{example}
\begin{solution}
配方得 $f(x)=(x-2)^{2}-1$，故对称轴为 $x=2$，顶点 $(2,-1)$。
\end{solution}

\begin{trapbox}{顶点横坐标的符号}
配方时括号内是 $x-1$ 而不是 $x+1$：$y=(x-1)^{2}-1$ 的对称轴是 $x=1$。
若误写为 $x=-1$，是把一次项系数的符号弄反了。
\end{trapbox}
